\documentclass[12pt, a4paper]{article}

% --- Packages ---
\usepackage[utf8]{inputenc}
\usepackage[T1]{fontenc}
\usepackage{lmodern}
\usepackage[margin=2.5cm]{geometry}
\usepackage{setspace}
\doublespacing

\usepackage{amsmath}
\usepackage{booktabs}       % professional tables
\usepackage{graphicx}       % figures
\usepackage{float}          % [H] placement
\usepackage{caption}
\usepackage{subcaption}     % subfigures
\usepackage{hyperref}
\usepackage{url}
\usepackage{natbib}         % author-year citations; swap for \usepackage[numbers]{natbib} if journal wants numbered
\usepackage{xcolor}
\usepackage{siunitx}        % consistent number/unit formatting e.g. \SI{45.5}{\EUR\per\MWh}
\usepackage{longtable}      % tables that span pages
\usepackage{multirow}
\usepackage{lineno}         % line numbers for review; remove for final
\linenumbers

% --- Metadata ---
\title{Levelized Cost of Heat Comparison for Electric and Natural Gas Pathways\\
       at EU E-PRTR Industrial Facilities:\\
       The Case for Off-Grid Solar with Thermal Heat Storage}

\author{
  Author Name\textsuperscript{1}\thanks{Corresponding author. \texttt{author@institution.edu}} \\
  \textsuperscript{1}Department, Institution, City, Country
  % Add co-authors as needed:
  % \and Co-Author\textsuperscript{2} \\
  % \textsuperscript{2}Department, Institution
}

\date{\today}

% =====================================================================
\begin{document}
\maketitle

% =====================================================================
\begin{abstract}
% TODO: Write last. ~200 words.
% Key points to hit:
%   (1) Motivation: European industrial process heat represents a major
%       decarbonization challenge.
%   (2) Gap: most analyses use grid electricity prices, obscuring the
%       economics of off-grid renewables.
%   (3) Methods: LCOH comparison across three pathways (natural gas boiler,
%       thermal heat battery, heat pump) for 110 EU E-PRTR/LCP facilities.
%   (4) Key finding: off-grid solar prices shift heat battery LCOH from
%       183.2 to 45.5 EUR/MWh, making it cheaper than natural gas at all
%       110 sites versus none with grid prices.
%   (5) Feasibility: 75 of 110 sites have sufficient land for on-site solar;
%       the main barrier for the remaining 35 is high heat demand.
%   (6) Policy implication: [TODO]

\noindent\textbf{Keywords:} industrial decarbonization; levelized cost of heat; off-grid solar; thermal heat battery; heat pump; E-PRTR; European industry
\end{abstract}

% =====================================================================
\section{Introduction}
\label{sec:intro}

% TODO: 3--5 paragraphs following this structure:
%
% Para 1 — Hook / motivation
%   Decarbonizing industrial process heat is one of the harder segments
%   of the energy transition. Cite Rosenow et al. and broader IEA/IPCC
%   context on hard-to-abate industrial heat.
%
% Para 2 — Current landscape
%   Heat pumps and electrification are increasingly discussed as pathways,
%   but cost competitiveness depends heavily on electricity prices.
%   Grid electricity in Europe (139--218 EUR/MWh) makes electrification
%   expensive relative to gas for most large consumers.
%
% Para 3 — Gap in literature
%   Off-grid generation (utility-scale solar + storage) as a direct
%   electricity source for industrial heat has received limited attention
%   at the facility level. Prior work has focused on national/sectoral
%   aggregates and grid pricing scenarios.
%
% Para 4 — This paper
%   We compare LCOH across three pathways — natural gas boiler, thermal
%   heat battery (solar + storage), and heat pump — for 110 large
%   combustion plant sites linked to E-PRTR facilities across Europe,
%   using both grid and off-grid electricity prices. We further assess
%   land feasibility for on-site solar generation at each facility.
%
% Para 5 — Roadmap sentence
%   Section~\ref{sec:data} describes data sources; Section~\ref{sec:methods}
%   the LCOH framework and feasibility methodology; Section~\ref{sec:results}
%   presents results; Section~\ref{sec:discussion} discusses implications
%   and limitations; Section~\ref{sec:conclusion} concludes.

% =====================================================================
\section{Data}
\label{sec:data}

\subsection{Facility Heat Demand (E-PRTR / LCP)}
% TODO:
%   - E-PRTR provides facility-level CO2 and pollutant data; LCP dataset
%     provides fuel input for combustion plants >50 MW rated thermal input.
%   - Describe the linking procedure between E-PRTR and LCP records.
%   - 110 matched facilities after excluding Activity 1c (power generation).
%   - Epsilon conversion factors: weighted-average fuel-to-heat efficiency
%     by sector and technology mix (cite eprtr_epsilon_reference).
%   - Caveat: LCP 50 MW threshold means heat demand estimates may
%     understate true sectoral totals.

\subsection{Natural Gas Prices}
% TODO:
%   - Eurostat dataset NRG\_PC\_203, Band I3, 2024.
%   - Range: 50--81 EUR/MWh LHV.
%   - LHV adjustment factor of 1.1098 applied to convert from GHV.
%   - Caveat: national average prices; may slightly overstate costs for
%     the largest consumers.

\subsection{Grid Electricity Prices}
% TODO:
%   - Eurostat dataset NRG\_PC\_205, industrial band, 2024.
%   - Range: 139--218 EUR/MWh.
%   - Same national-average caveat as gas prices.

\subsection{Off-Grid Electricity Prices (LCOE)}
% TODO:
%   - Calculated from Bloomberg NEF (BNEF) 2025 cost data: capex, annual
%     opex, and cost of capital for utility-scale PV (fixed-axis) and
%     battery storage, by country.
%   - 84 sites matched directly to BNEF country data; 26 proxied from
%     peer countries (by income level).
%   - Solar capacity factors from PVGIS ERA5 API (lat/lon, fixed-axis,
%     most recent complete year: 2023).
%   - EUR/USD conversion: ECB 2024 annual average = 0.921.

\subsection{Technology Cost Assumptions}

% TODO: brief prose description of sources before the table.

\begin{table}[H]
  \centering
  \caption{Technology capital and operating cost assumptions.}
  \label{tab:tech_costs}
  \begin{tabular}{llll}
    \toprule
    \textbf{Technology} & \textbf{Capex} & \textbf{Opex} & \textbf{Source} \\
    \midrule
    Gas boiler      & 75 \euro/kW$_\text{th}$                  & 1.5\% of capex & [TODO] \\
    Thermal heat battery & 150 \euro/kW$_\text{th}$, 8-h storage & 0.5\% of capex & NREL 2025 \citep{wikoff2025} \\
    Heat pump       & 1{,}200 \euro/kW$_\text{th}$              & 2.5\% of capex & IRENA 2022 \citep{irena2022hp} \\
    \bottomrule
  \end{tabular}
\end{table}

% =====================================================================
\section{Methods}
\label{sec:methods}

\subsection{Levelized Cost of Heat (LCOH)}
% TODO:
%   Write out the LCOH formula explicitly. Something like:
%
%   \begin{equation}
%     \text{LCOH} = \frac{\text{CRF} \cdot \text{Capex} + \text{Opex}}
%                        {E_\text{heat} \cdot \text{CF}_\text{plant}}
%                   + \frac{P_\text{fuel}}{\eta}
%                   + \tau \cdot \text{EF}
%   \end{equation}
%
%   where CRF is the capital recovery factor (8\%, 20-year lifetime),
%   $E_\text{heat}$ is annual heat demand, CF$_\text{plant}$ = 0.85
%   (industrial process capacity factor), $P_\text{fuel}$ is the energy
%   price (EUR/MWh), $\eta$ is conversion efficiency (boiler: 0.9),
%   $\tau$ is the EU ETS carbon price (60 EUR/tCO$_2$ in 2024), and
%   EF is the emission factor (natural gas: 0.202 tCO$_2$/MWh LHV).
%
%   Note: worked example in the draft (gas price = 30.63 EUR/MWh) is a
%   good sanity-check box or appendix item.

\subsection{Land Feasibility Assessment}
% TODO:
%   - Land cover: grassland, shrubland, cropland, bare land (exclude
%     protected areas, slopes >5°).
%   - 10 km buffer around each facility.
%   - Overlap rule: shared buffer area allocated proportionally to
%     heat demand.
%   - Feasibility criterion: suitable land area $\geq$ land required
%     to meet site heat demand given local solar capacity factor.
%   - Land required derived from solar capacity factor (PVGIS) and
%     assumed panel density [TODO: state W/m² assumption].

% =====================================================================
\section{Results}
\label{sec:results}

\subsection{LCOH Comparison: Grid vs.\ Off-Grid Electricity}

% TODO: 1--2 paragraphs interpreting Table~\ref{tab:lcoh}.

\begin{table}[H]
  \centering
  \caption{Mean LCOH by pathway and electricity price scenario (EUR/MWh$_\text{th}$).}
  \label{tab:lcoh}
  \begin{tabular}{lcc}
    \toprule
    \textbf{Pathway}       & \textbf{Off-grid price} & \textbf{Grid price} \\
    \midrule
    Natural gas boiler     & 83.8  & 83.8  \\
    Thermal heat battery   & 45.5  & 183.2 \\
    Heat pump              & 70.8  & 79.3  \\
    \bottomrule
  \end{tabular}
\end{table}

% TODO: note that at grid prices 99/110 sites have cheapest LCOH from
% natural gas, 11 from heat pump; at off-grid prices all 110 sites are
% cheapest with heat battery.

\subsection{Site-Level Feasibility}

% TODO: describe the 75 feasible / 35 infeasible split and Table~\ref{tab:feasibility}.

\begin{table}[H]
  \centering
  \caption{Characteristics of feasible vs.\ infeasible sites.}
  \label{tab:feasibility}
  \begin{tabular}{lcc}
    \toprule
    & \textbf{Infeasible (35)} & \textbf{Feasible (75)} \\
    \midrule
    Mean heat demand (MWh$_\text{th}$/yr) & 1{,}092{,}000 & 300{,}000 \\
    Mean land required (km²)              & 74.3          & 22.4      \\
    Mean suitable land in buffer (km²)    & 26.2          & 111.5     \\
    Mean coverage ratio                   & 0.32          & 194$\times$ (median: 7.7$\times$) \\
    Mean suitability fraction of buffer   & 8.4\%         & 35.5\%    \\
    \bottomrule
  \end{tabular}
\end{table}

% TODO: discuss drivers of infeasibility:
%   - High heat demand is the primary barrier.
%   - Coastal/port locations (refineries, petrochemicals) have limited
%     suitable land adjacent.
%   - Low-temperature industries (sugar, pulp \& paper, glass,
%     waste-to-energy) most feasible: rural locations, lower heat demand.

\subsection{Geographic and Sectoral Patterns}
% TODO: describe map figure here. Placeholder below.

\begin{figure}[H]
  \centering
  % \includegraphics[width=\textwidth]{figures/facility_map.pdf}
  \fbox{\parbox{0.9\textwidth}{\centering [TODO: map of facilities coloured by cheapest LCOH pathway and feasibility]}}
  \caption{Map of E-PRTR/LCP facilities showing least-cost heat pathway
           (off-grid scenario) and solar feasibility.}
  \label{fig:map}
\end{figure}

% =====================================================================
\section{Discussion}
\label{sec:discussion}

\subsection{Why Off-Grid Pricing Changes the Economics}
% TODO: explain the mechanism — off-grid LCOE avoids grid charges,
% network tariffs, and taxes that inflate industrial electricity prices.
% Quantify the price gap (grid: 139--218 vs. off-grid: [TODO] EUR/MWh).

\subsection{Limitations}
% TODO — be explicit; journals appreciate this:
\begin{itemize}
  \item \textbf{National-average energy prices:} Gas and grid electricity
        prices from Eurostat Band I3 are national averages; facility-level
        prices may differ, particularly for the largest consumers.
  \item \textbf{LCP 50 MW threshold:} Only plants with rated thermal input
        $>$50 MW are included, likely understating true sectoral heat demand.
  \item \textbf{Solar capacity factor vintage:} PVGIS data from 2023; may
        not reflect long-run averages at all sites.
  \item \textbf{No grid backup costs:} The off-grid scenario assumes
        full self-supply; grid backup or curtailment costs are not modelled.
  \item \textbf{Heat demand derivation:} Epsilon factors assume a fixed
        technology mix per sector; actual facility mix may differ.
  \item \textbf{Static cost assumptions:} BNEF 2025 costs used throughout;
        technology costs are falling and results will evolve.
\end{itemize}

\subsection{Policy Implications}
% TODO: where should industrial decarbonization incentives be concentrated?
% Which sectors / geographies are lowest-hanging fruit?

% =====================================================================
\section{Conclusion}
\label{sec:conclusion}

% TODO: 2--3 paragraphs. No new information.
%   Para 1 — restate the question and core finding.
%   Para 2 — feasibility result and what drives it.
%   Para 3 — policy takeaway and future work.

% =====================================================================
\section*{Data Availability}
% TODO: state where data and code will be deposited (e.g., Zenodo, GitHub).

\section*{Acknowledgements}
% TODO

\section*{Declaration of Interests}
The authors declare no competing interests.

% =====================================================================
\bibliographystyle{apalike}   % or use a journal-specific .bst file
\bibliography{references}     % point to your .bib file

% =====================================================================
\appendix
\section{LCOH Calculation: Worked Example}
\label{app:lcoh_example}

% TODO: reproduce the worked example from the draft here as a sanity check.
%   Gas price = 30.63 EUR/MWh LHV
%   Energy component: (1/0.9) * 30.63 = 34.03 EUR/MWh_th
%   Carbon component: 0.202 tCO2/MWh * 60 EUR/tCO2 = 12.12 EUR/MWh_th
%   Fixed component: CRF * capex / (CF_plant * 8760) + opex ... [complete]

\section{Epsilon Conversion Factors by Sector}
\label{app:epsilon}
% TODO: reproduce the epsilon reference table from eprtr_epsilon_reference.csv.

\section{Country Proxying for BNEF Data}
\label{app:proxying}
% TODO: list the 26 proxied countries and their peer-country matches.

\end{document}
